\paragraph{RQ1: Construction and capacity.}
How can JHQ be built on a GPU, and what cost, resident capacity and build peak result (\cref{sec:7-7})?

\paragraph{RQ2: Correct and efficient execution.}
Which execution choices speed up search while preserving score/tie contracts (\cref{sec:7-3,sec:7-4,sec:7-6})?

\paragraph{RQ3: Budget quality and total cost.}
Can unlabelled calibration preserve held-out quality and repay its cost (\cref{sec:7-5})? The GT-labelled oracle is an evaluation control.

\paragraph{RQ4: Competitive operating regions.}
Under host-to-host timing, where does JHQ lead at matched recall?
The answer uses independently tuned E1 frontiers and explicit memory
definitions (\cref{sec:7-2}). Int8 graph competitors, high-recall reversals
and absent coverage are part of the answer.

Chapter 3 turns these questions into controlled comparisons. Together they
separate what was built, whether its execution is correct, what effort a
workload requires, and whether the resulting system has useful operating
regions.
